\documentclass[10pt,a4paper]{amsart}
\usepackage[margin=19mm]{geometry}
\usepackage{amsmath,amssymb,mathtools}
\usepackage[scr=rsfs,cal=euler]{mathalfa}
\usepackage{xfrac}
\usepackage[unicode,hidelinks,bookmarks=false]{hyperref}
\newcommand{\ie}{i.e.\ }
\newcommand{\cf}{cf.\ }
\newcommand{\resp}{respectively\ }
\newcommand{\Subsection}{\subsection}
\providecommand{\nobreakdash}{\nobreak\hbox{-}}
\setlength{\emergencystretch}{2em}
\multlinegap0pt
\usepackage{tikz}
\usepackage[shortlabels]{enumitem}
\usepackage{float}
\usepackage{amssymb}
\newtheorem{lemma}{Lemma}
\title{Limitations of deducing measures of limsup sets from measures of finite intersections}
\author{Charlie Wilson}
\date{Source: arXiv:2502.02416v3 (CC BY 4.0)}
\begin{document}
\maketitle

\noindent\emph{Source and licence.} Limitations of deducing measures of limsup sets from measures of finite intersections. Version arXiv:2502.02416v3 (CC BY 4.0), distributed by arXiv under the Creative Commons Attribution 4.0 International licence, http://creativecommons.org/licenses/by/4.0/, as recorded in the arXiv licence metadata for this exact version and shown on the abstract page https://arxiv.org/abs/2502.02416v3. Full licence notice: \texttt{licenses/arxiv-2502.02416-CC-BY-4.0.txt}.\par\smallskip

\noindent\textit{Editorial scope.} Counted unit: the article's lemma lemma (source0.tex lines 303--315). The article's complete original proof of it is delivered with it (source0.tex lines 322--465, one \texttt{proof} environment of 144 lines, which the article places away from the statement). No mathematics is written, repaired or re-derived: the statement and the proof are the article's own lines, in the article's own notation, and no retained line is substituted or re-flowed.  No further source-local context is needed: every cross-reference of every retained line resolves inside the two delivered spans.  Nothing was omitted from any retained span, so the package carries no omission record.  No retained line contains a citation, so the wrapper carries no bibliography and no bibliography entry.  The retained lines contain the article's own diagram code, which the wrapper typesets with the standard tikz packages; no image file is delivered and the package declares no asset dependency.  Editorial formatting only, in addition: an AMS \texttt{amsart} preamble that restates the article's own declarations and macros used by the retained lines, namely the environments \texttt{lemma} on separate counters, together with the standard packages those lines need.  The retained environments print in this document as Lemma 1 in order of appearance, whereas the article prints the same items with the numbering of its own full document.  The complete native source of the article as distributed in the arXiv e-print for this exact version is bundled unchanged as \texttt{sources/172733/source0.tex}.
\begin{lemma}\label{prelim equality lemma}
    For any $m \in \mathbb{N}$, there exist two finite collections  $C_1,\allowbreak C_2,\allowbreak \dots,\allowbreak C_{m+1}$ and
$D_1,\allowbreak D_2,\allowbreak \dots,\allowbreak D_{m+1}$ of subsets of $[0,1]$ such that the following hold: 
    \begin{enumerate}[label=(\roman*)]
    \item $\mu (C_i) = \mu (D_i) $ for all $1\leq i \leq m+1$ ,
     \item For any sequence of indices $1 \leq i_1,i_2,...,i_l \leq m+1$ with $1 \leq l \leq m$, we have \[\mu(C_{i_1} \cap ... \cap C_{i_l}) = \mu(D_{i_1} \cap ... \cap D_{i_l}),\]

     \item $1=\mu \left( \bigcup_{i=1}^{m+1} D_i \right) = \mu \left( \bigcup_{i=1}^{m+1} C_i \right)+\frac{1}{2^m}, and$     
     \item  for each $1 \leq i \leq m+1$, $C_i$ and $D_i$ are unions of dyadic intervals of length $2^{-m}$.
     \end{enumerate}


\end{lemma}

\begin{proof}[Proof of Lemma \ref{prelim equality lemma}]
    We start by considering the collection $\{D_1,D_2,...,D_{m+1}\}$.
    We list all subsets of \mbox{$\{D_1,\allowbreak D_2,\allowbreak \dots,\allowbreak D_{m+1}\}$} consisting of an odd number of $D_i$ and call this collection $D^{\textrm{odd}}$. Similarly, we define $D^{\textrm{even}}$ to be the collection of subsets of $\{D_1,D_2,\dots,D_{m+1}\}$ consisting of an even number of elements. For each of the subsets in $D^{\textrm{odd}}$, we assign a distinct (not assigned to any of the other subsets) interval of size $2^{-m}$ from $\mathcal{D}_m$ and then add this interval to all the $D_i$ in this subset. We note that there are $2^m$ nonempty subsets of odd size and also $2^m$ dyadic intervals; therefore, this procedure is well defined and ensures that each interval is contained in at least one of the $D_i$. This ensures that:
    \[\mu \left( \bigcup_{i=1}^{m+1}D_i \right) =1.\]
    To aid the reader, for the case $m=3$ we include an example of the correspondence between the collection $D^{\textrm{odd}}$ to the dyadic intervals $\mathcal{D}_3$ in Table $\ref{correspondence table}$. Furthermore, we include the analogous examples of the sets $D_i$ in Figure \ref{Di set example m=3}.

\begin{table}[h]
        \centering \caption{An example of the correspondence between  between $D^{\textrm{odd}}$ and the dyadic intervals $\mathcal{D}_m$ in the case $m=3$. }
\vspace{10pt}
{\renewcommand{\arraystretch}{2}
    \begin{tabular}{|c|c|c|c|c|}
\hline
Collection of $D_i$ & $\{D_1 \}$ & $\{D_2\}$ & $\{D_1,D_3,D_4\}$  & $\{D_2,D_3,D_4\}$  \\
\hline
Associated interval & $ \left[ \frac{3}{8}, \frac{4}{8} \right] $ & $ \left[ \frac{5}{8}, \frac{6}{8} \right] $ & $ \left[ \frac{2}{8}, \frac{3}{8} \right] $ & $ \left[ \frac{4}{8}, \frac{5}{8} \right] $ \\
\hline
\hline
Collection of $D_i$ & $\{D_3 \}$ & $\{D_4\}$ & $\{D_1,D_2,D_3\}$  & $\{D_1,D_2,D_4\}$  \\
\hline
Associated interval & $ \left[ \frac{6}{8}, \frac{7}{8} \right] $ & $ \left[ \frac{7}{8}, \frac{8}{8} \right] $ & $ \left[ \frac{0}{8}, \frac{1}{8} \right] $ & $ \left[ \frac{1}{8}, \frac{2}{8} \right] $ \\
\hline
    \end{tabular}}
    \label{correspondence table}
\end{table}

\begin{figure}[!ht]
    \caption{An example of the sets $D_i$ in the case $m=3$.}
    \label{leb construction Di}
\begin{center}
    \begin{tikzpicture}
    \draw[very thick] (0,5) -- (12,5);
    \draw (0,4) -- (12/8*4,4);
    \draw (0,3) --(2*12/8,3);
    \draw (4*12/8,3) --(6*12/8,3);
    \draw (0,2) -- (12/8,2);
    \draw (2*12/8,2) -- (3*12/8,2);
    \draw (4*12/8,2) -- (5*12/8,2);
    \draw (6*12/8,2) -- (7*12/8,2);
    \draw (12/8,1) -- (2*12/8,1);
    \draw (2*12/8,1) -- (3*12/8,1);
    \draw (4*12/8,1) -- (5*12/8,1);
    \draw (7*12/8,1) -- (8*12/8,1);
    \node at (-1/2,5) {[0,1]};
    \node at (-1/2,4) {$D_1$};
    \node at (-1/2,3) {$D_2$};
    \node at (-1/2,2) {$D_3$};
    \node at (-1/2,1) {$D_4$};
    \node at (0,1/2) {0};
    \node at (12/8,1/2) {$\frac{1}{8}$};
    \node at (24/8,1/2) {$\frac{2}{8}$};
    \node at (3*12/8,1/2) {$\frac{3}{8}$};
    \node at (4*12/8,1/2) {$\frac{4}{8}$};
    \node at (5*12/8,1/2) {$\frac{5}{8}$};
    \node at (6*12/8,1/2) {$\frac{6}{8}$};
    \node at (7*12/8,1/2) {$\frac{7}{8}$};
    \node at (12,1/2) {1};
    \end{tikzpicture}
\end{center}
\label{Di set example m=3}
\end{figure}
    
    We construct the $C_i$ in a very similar vein. 
    We consider the collection $\{ C_1,...,C_{m+1} \}$ and define $C^{\textrm{odd}}$ and $C^{\textrm{even}}$ analogously to $D^{\textrm{odd}}$ and $D^{\textrm{even}}$ above. For each subset in $C^{\textrm{even}}$, we assign a distinct (not assigned to any of the other subsets) interval of size $2^{-m}$ from $\mathcal{D}_m$ and then add this interval to all the $C_i$ in this subset. However, now there are $2^m-1$ non-empty subsets of even size, so there will be an interval of size $2^{-m}$ which will not be in any of the $C_i$. Thus, it follows that:
     \[\mu \left( \bigcup_{i=1}^{m+1} C_i \right) =1-\frac{1}{2^m}.\]


\begin{figure}[h!]
    \caption{An example of the sets $C_i$ for $m=3$.}
    \label{leb construction Ci}
\begin{center}
    \begin{tikzpicture}
    \draw[very thick] (0,5) -- (12,5);
    \draw (0,4) -- (12/8*4,4);
    \draw (0,3) --(12/8*2,3);
    \draw (4*12/8,3) --(6*12/8,3);
    \draw (0,2) -- (12/8,2);
    \draw (2*12/8,2) -- (3*12/8,2);
    \draw (4*12/8,2) -- (5*12/8,2);
    \draw (6*12/8,2) -- (7*12/8,2);
    \draw (0,1) -- (12/8,1);
    \draw (3*12/8,1) -- (4*12/8,1);
    \draw (5*12/8,1) -- (6*12/8,1);
    \draw (6*12/8,1) -- (7*12/8,1);
    \node at (-1/2,5) {[0,1]};
    \node at (-1/2,4) {$C_1$};
    \node at (-1/2,3) {$C_2$};
    \node at (-1/2,2) {$C_3$};
    \node at (-1/2,1) {$C_4$};
    \node at (0,1/2) {0};
    \node at (12/8,1/2) {$\frac{1}{8}$};
    \node at (24/8,1/2) {$\frac{2}{8}$};
    \node at (3*12/8,1/2) {$\frac{3}{8}$};
    \node at (4*12/8,1/2) {$\frac{4}{8}$};
    \node at (5*12/8,1/2) {$\frac{5}{8}$};
    \node at (6*12/8,1/2) {$\frac{6}{8}$};
    \node at (7*12/8,1/2) {$\frac{7}{8}$};
    \node at (12,1/2) {1};
    \end{tikzpicture}
\end{center}
\end{figure}




 Note that conditions $(iii)$ and $(iv)$ from the lemma statement are satisfied by the construction of the sets $C_i$ and $D_i$. Also note that condition $(i)$ follows from condition $(ii)$. Thus, it remains to show that condition $(ii)$ is satisfied. By the symmetry of the construction of the sets $C_i$ and $D_i$, it suffices to show that 
 \[\mu(D_1 \cap ... \cap  D_l) = \mu(C_1 \cap ... \cap  C_l) \qquad \text{for } 1 \leq l \leq m.\]

For a subset $X$ of $[0,1]$, we write $X^c$ to denote its complement, i.e. $X^c=[0,1]\setminus X$. For convenience, it will sometimes be helpful for us to write $X^1$ in place of $X$. We can write $D_1 \cap D_2 \cap \dots \cap D_l$ as the following disjoint union:
 \begin{align*}
     D_1 \cap D_2 \cap \dots \cap D_l &= \bigsqcup_{\substack{\sigma_i \in \{1,c\} \\ l+1 \leq i \leq m+1}} D_1 \cap D_2 \cap ...\cap D_l \cap D_{l+1}^{\sigma_{l+1}} \cap ... \cap D_{m+1}^{\sigma_{m+1}}. 
 \end{align*}
It follows that 
    \begin{align}\label{D decompose}
        \mu(D_1 \cap...\cap D_l)=\sum_{\substack{\sigma_i \in \{1,c\} \\ l+1 \leq i \leq m+1}} \mu\left(D_1\cap...\cap D_l \cap D_{l+1}^{\sigma_{l+1}} \cap ...\cap D_{m+1}^{\sigma_{m+1}} \right).
    \end{align}
We see from the construction of the $D_i$ that for $1 \leq l \leq m$ we have
\begin{equation}\label{D measure}
\mu\left(D_1\cap...\cap D_l \cap D_{l+1}^{\sigma_{l+1}} \cap ... \cap D_{m+1}^{\sigma_{m+1}}\right)=
    \begin{cases}
        \frac{1}{2^m}  & \text{if } ~ l+ \#\{i:\sigma_i=1\} ~\text{is odd,}\\
        0 & \text{else. }
    \end{cases}
\end{equation}
This is because in order for the left-hand side of (\ref{D measure}) to be non-zero we need the number of non-complement sets being included in the intersection to be odd and precisely one interval of length $2^{-m}$ lies in this intersection from the construction.
In a similar vein we can write:
\begin{align*}
     C_1 \cap C_2 \cap \dots \cap C_l &= \bigsqcup_{\substack{\sigma_i \in \{1,c\} \\ l+1 \leq i \leq m+1}} C_1 \cap C_2 \cap ...\cap C_l \cap C_{l+1}^{\sigma_{l+1}} \cap ... \cap C_{m+1}^{\sigma_{m+1}}. 
 \end{align*} 
 Thus,
\begin{align}\label{C decompose}
    \mu\left(C_1 \cap...\cap C_l\right)=\sum_{\substack{\sigma_i \in \{1,c\} \\ l+1 \leq i \leq m+1}} \mu\left(C_1\cap...\cap C_l \cap C_{l+1}^{\sigma_{l+1}} \cap ...\cap C_{m+1}^{\sigma_{m+1}}\right).
\end{align}
We see from the construction of the $C_i$ that for $1 \leq l \leq m$ we have
\begin{equation}\label{C measure}
\mu\left(C_1\cap...\cap C_l \cap C_{l+1}^{\sigma_{l+1}} \cap...\cap C_{m+1}^{\sigma_{m+1}} \right)=
    \begin{cases}
        \frac{1}{2^m}  & \text{if } ~ l+ \#\{i:\sigma_i=1\} ~\text{is even,}\\
        0 & \text{else. }
    \end{cases}
\end{equation}
By similar reasoning to that of $D_i$, in order for the left-hand side of (\ref{C measure}) to be non-zero we need the number of non-complement sets being included in the intersection to be even and precisely one interval of length $2^{-m}$ lies in this intersection from the construction.

 For the set $\{ \sigma_{l+1},...,\sigma_{m+1} \}$, the number of subsets of odd size is equal to the number of subsets of even size, and hence by combining the results $(\ref{D measure})$ and $(\ref{C measure})$ with the decompositions $(\ref{C decompose})$ and $(\ref{D decompose})$, $\mu(C_1\cap...\cap C_l)=\mu(D_1\cap...\cap D_l)$ for $1 \leq l \leq m$ as required.
\end{proof}

\end{document}
